% ============================================================
%  THESIS METADATA
%  Edit the values in this file for your thesis. Everything
%  else (preamble.tex, titlepages.tex) is reusable template
%  code and should not normally need to change.
% ============================================================

% --- People ---
\newcommand{\Name}{Andrei-Silviu Voicu}
\newcommand{\Advisor}{Conf. dr. ing. Cornel Popescu}

% --- Title (Romanian / English) ---
\newcommand{\ProjectTitleRO}{Platformă educațională pentru învățarea limbii engleze}
\newcommand{\ProjectTitleEN}{Educational platform for learning English}

% --- Date shown on the cover pages ---
\newcommand{\MonthRO}{IULIE}
\newcommand{\MonthEN}{JULY}
\newcommand{\Year}{2026}

% --- Abstracts ---
% Max. ~200 words each; together both versions should fit on one page.
\newcommand{\AbstractRO}{Lucrarea prezintă o platformă web pentru studiul limbii engleze. Ideea de la care s-a plecat este că majoritatea platformelor existente acoperă fie conținutul didactic, fie sesiunile live cu profesori, rareori pe ambele, ceea ce îl obligă pe utilizator să combine mai multe servicii ca să își acopere nevoile de bază. Platforma include cursuri pe trei niveluri de dificultate, exerciții și teste cronometrate, iar utilizatorii pot rezerva sesiuni cu profesori verificați pe platformă, fie individual, fie în grup, accesul la sesiuni funcționând pe baza unui sistem de credite. Profesorii nu pot preda imediat după înregistrare, ci trec printr-un proces de aprobare, iar cursurile create de aceștia sunt revizuite înainte de a apărea în catalog. Progresul fiecărui student este urmărit atât la nivel de lecție, cât și la nivel de quiz, rezultatul final fiind o aplicație care acoperă ciclul întreg de învățare, de la înregistrarea contului până la susținerea sesiunilor video.}

\newcommand{\AbstractEN}{This paper presents a web platform for learning English. The starting point was a fairly common problem: most existing platforms cover either study materials or live sessions with instructors, rarely both, which forces users to combine multiple services just to cover their basic needs. The platform includes courses across three difficulty levels, interactive exercises and timed quizzes, and users can book sessions with verified instructors, either one-on-one or in a group, with access working through a credit system tied to a subscription. Instructors cannot start teaching right after signing up, they go through an approval process, and any courses they create are reviewed before appearing in the catalog. Student progress is tracked at both the lesson and quiz level, with the end result being a functional application that covers the full learning cycle, from account registration to completing video sessions with an instructor.}
